\documentclass[10pt,a4paper]{amsart}
\usepackage[margin=19mm]{geometry}
\usepackage{amsmath,amssymb,mathtools}
\usepackage[scr=rsfs,cal=euler]{mathalfa}
\usepackage{xfrac}
\usepackage[unicode,hidelinks,bookmarks=false]{hyperref}
\newcommand{\ie}{i.e.\ }
\newcommand{\cf}{cf.\ }
\newcommand{\resp}{respectively\ }
\newcommand{\Subsection}{\subsection}
\providecommand{\nobreakdash}{\nobreak\hbox{-}}
\setlength{\emergencystretch}{2em}
\multlinegap0pt
\usepackage{bm}
\usepackage{bbm}
\usepackage{dsfont}
\usepackage{xspace}
\usepackage{cancel}
\usepackage{mathtools}
\usepackage{amsthm}
\makeatletter
\@ifundefined{theorem}{\newtheorem{theorem}{Theorem}}{}
\@ifundefined{proposition}{\newtheorem{proposition}[theorem]{Proposition}}{}
\makeatother
\def\al{\alpha}
\def\be{\beta}
\def\g{\mathfrak{g}}
\def\lmd{\lambda}
\def\vf{\varphi}
\newcommand{\C}{\mathbb C}
\newcommand{\Z}{\mathbb{Z}}
\newcommand{\elei}[2]{I_{#1}^{(#2)}}
\newcommand{\ov}[1]{\overline{#1}}
\title[Non-weight modules over gap-$p$ Virasoro algebras]{Non-weight modules over gap-$p$ Virasoro algebras}
\author{Chengkang Xu, Fulin Chen, Shaobin Tan}
\date{Source: arXiv:2505.15273v1 (CC BY 4.0)}
\begin{document}
\maketitle

\noindent\emph{Source and licence.} Non-weight modules over gap-$p$ Virasoro algebras. Version arXiv:2505.15273v1 (CC BY 4.0), distributed under the Creative Commons Attribution 4.0 International licence, as recorded in the arXiv licence metadata for this exact version and shown on the abstract page https://arxiv.org/abs/2505.15273v1. Full rights notice: licenses/arxiv-2505.15273-CC-BY-4.0.txt.\par\smallskip

\noindent\textit{Editorial scope.} Counted unit: the Proposition whose source label is \texttt{prop5.2} in the article (source0.tex lines 941--945, source label \texttt{prop5.2}), delivered together with the article's own complete original proof of it (source0.tex lines 946--1003, one \texttt{proof} environment of 58 lines). No mathematics is written, repaired or re-derived: statement and proof are the article's own text line for line. Required context retained uncounted: none. Every definition, display and label of those lines is retained unchanged. Omitted from this package and not redistributed: 1--940, 1004--1624. Nothing inside a retained excerpt is omitted or altered: every retained body is delivered exactly as the article's lines are written, so no omission receipt of any kind accompanies this package. Citations: the retained excerpts carry no citation command at all, so no bibliography entry of the article is transcribed and no reference is redistributed. Reference adaptation: none was needed and none was made; every reference inside the delivered unit resolves inside it, so no unresolved reference or citation is produced. Printed slips kept literal: no printed word, symbol, hypothesis or number of the counted statement or of its proof was altered. Layout and class: the article's own class is \texttt{amsart} with packages including amssymb, dsfont, eucal, amsmath, amscd, color, multicol, xy, graphicx, colordvi, xspace, txfonts, lscape, tikz; the wrapper is the lane's pinned \texttt{amsart} editorial wrapper at 10pt on A4, which restates the theorem-like environments the retained text needs and adds the article's own macro definitions that the retained text needs (\texttt{\textbackslash C}, \texttt{\textbackslash Z}, \texttt{\textbackslash al}, \texttt{\textbackslash be}, \texttt{\textbackslash elei}, \texttt{\textbackslash g}, \texttt{\textbackslash lmd}, \texttt{\textbackslash ov}, \texttt{\textbackslash vf}), with extra wrapper packages (bm, bbm, dsfont, xspace, cancel, mathtools). The delivered numbering is the unit's own. Assets: no retained span contains a figure environment, a graphics-inclusion command, a PDF-page inclusion, a table or a TikZ picture; no image file is copied into this package, the package declares no asset dependency and no image file is redistributed with it. Licence: version arXiv:2505.15273v1 of this article is distributed under the Creative Commons Attribution 4.0 International licence, as recorded in the arXiv licence metadata for this exact version and shown on the abstract page https://arxiv.org/abs/2505.15273v1. A full rights notice is delivered at licenses/arxiv-2505.15273-CC-BY-4.0.txt. Characters: every retained excerpt is pure ASCII, so the delivered engine, XeLaTeX with its default Latin Modern text font, reports no missing character.
\begin{proposition}\label{prop5.2}
 Let $\lmd,\mu\in\C^\times, \ov\al,\ov\beta\in\C^p$ and $V,W$ be irreducible restricted modules over $\g$.
 Then  $\Omega(\lmd,\ov\al)\otimes V\cong\Omega(\mu,\ov\beta)\otimes W$
 if and only if $\lmd=\mu, \ov\al=\ov\beta$, and $V\cong W$.
\end{proposition}

\begin{proof}
 It suffices to prove the ``only if" part, and for that we take a $\g$-module isomorphism
 \[\vf:\Omega(\lmd,\ov\al)\otimes V\longrightarrow\Omega(\mu,\ov\beta)\otimes W.\]
 For any $v\in V$, write
 \begin{equation}\label{eq:exvf}
    \vf(1\otimes v)=\sum_{i=0}^st^i\otimes w_i
 \end{equation} with $w_i\in W,s\geq 0$ and $w_s\neq 0$.
 Take $N_v\in\Z_+$ such that $L_nv=L_n(w_i)=\elei njv=\elei njw_i=0$ for all $n\geq N_v,1\leq j\leq p-1,0\leq i\leq s$.

 We first show that $\lambda=\mu$.
 For any $m,n\geq N_v$, by combining \eqref{eq:exvf} with the equality
 \[\vf\left(\left(\lmd^{-m}L_m-\lmd^{-n}L_n\right)(1\otimes v)\right)
     =\left(\lmd^{-m}L_m-\lmd^{-n}L_n\right)\vf(1\otimes v),\]
     we get that
 \begin{equation}\label{eq5.1}
   \al_0(n-m)\sum_{i=0}^st^i\otimes w_i
   =\left(\frac\mu\lmd\right)^m\sum_{i=0}^s(t+\be_0m)(t+m)^i\otimes w_i
     -\left(\frac\mu\lmd\right)^n\sum_{i=0}^s(t+\be_0n)(t+n)^i\otimes w_i.
 \end{equation}
 Notice that the coefficient of $t^{s+1}\otimes w_s$ in the right hand side of \eqref{eq5.1} is
 $\left(\frac\mu\lmd\right)^m-\left(\frac\mu\lmd\right)^n$, while that in the left hand side is 0.
This implies that  $\lmd=\mu$.

Next we prove that $\ov\alpha=\ov\beta$.
Since  $\lmd=\mu$, \eqref{eq5.1} can be rewritten as
 \begin{equation}\label{eq5.2}
   \al_0(n-m)\sum_{i=0}^st^i\otimes w_i=\sum_{i=0}^s(t+\be_0m)(t+m)^i\otimes w_i
          -\sum_{i=0}^s(t+\be_0n)(t+n)^i\otimes w_i.
 \end{equation}
By comparing  the coefficients of $t^s\otimes w_s$ in the both sides of \eqref{eq5.2}, we find that
 $\al_0(n-m)=(\be_0+s)(n-m)$ and so $\al_0=\be_0+s\geq \be_0$.
 Replacing $\vf$ by $\vf^{-1}$ in the above discussion, we also get $\be_0\geq \al_0$.
 This gives  $\be_0=\al_0$ and $s=0$. Particularly, we have  \[\vf(1\otimes v)=1\otimes w_0.\]
Furthermore, for every $n\geq N_v $ and $1\leq i\leq p-1$, we obtain
 $$\be_i\lmd^n(1\otimes w_0)=\elei ni(1\otimes w_0)=\elei ni\vf(1\otimes v)
      =\vf\left(\elei ni(1\otimes v)\right)=\al_i\lmd^n(1\otimes w_0).$$
This shows $\be_i=\al_i$, as desired.

Finally, we show that $V\cong W$.
Let us define a linear map $\eta:V\longrightarrow W$ by
 $$\vf(1\otimes w)=1\otimes \eta(w)\qquad\text{ for any }w\in V.$$
 It is clear that $\eta$ is injective.
 For every $n\in\Z$ and $1\leq i\leq p-1$,
it follows from the equality $\vf\left(\elei ni(1\otimes v)\right)=\elei ni\vf(1\otimes v)$ that
 $$\al_i\lmd^n\vf(1\otimes v)+\vf\left(1\otimes \elei ni v\right)=\elei ni(1\otimes\eta(v))
   =\al_i\lmd^n(1\otimes\eta(v))+1\otimes\elei ni\eta(v).$$
This implies that $1\otimes\eta\left(\elei ni v\right)=\vf\left(1\otimes \elei ni v\right)=1\otimes\elei ni\eta(v)$, and hence
 $\eta\left(\elei ni v\right)=\elei ni\eta(v)$.

 Similarly, the equality $\vf(L_n(1\otimes v))=L_n\vf(1\otimes v)$ for any $n\geq N_v$ implies that
 $$\vf(t\otimes v)=t\otimes \eta(v).$$
Additionally, by applying the equality  $\vf(L_n(1\otimes v))=L_n\vf(1\otimes v)$ for any $n\in \Z$  we have
 $\vf(1\otimes L_nv)=1\otimes L_n\eta(v)$. Thus
 $$\eta(L_nv)=L_n\eta(v)\qquad\text{ for any }n\in\Z.$$
 This proves that $\eta$ is a $\g$-module homomorphism.
 Since $\eta(V)\neq 0$ and $W$ is irreducible, $\eta$ is an isomorphism.
 We complete the proof.
\end{proof}

\end{document}
